\subsection{四元数代数的单性}

命题 3. —— 设 \(\alpha, \beta, \gamma\) 是 K 的元素，\(f\) 是型为 \((\alpha, \beta, \gamma)\) 的四元数代数。把它的凯莱迹与凯莱范数记作 \(\mathrm{T}_{\mathrm{F}}\) 与 \(\mathrm{N}_{\mathrm{F}}\)。下列性质等价：

\begin{enumerate}
\def\labelenumi{(\roman{enumi})}
\item
  代数 \(\mathbf{F}\) 是中心单的。
\item
  对每个非零元素 \(x\)（属于 \(\mathrm{F}\)），存在一个元素 \(y\)（在 \(\mathrm{F}\) 中），使得 \(\mathrm{T}_{\mathrm{F}}(xy) \neq 0\)。
\item
  我们有 \((4\alpha + \beta^{2})\gamma \neq 0\)。
\end{enumerate}

设这些性质成立。则对 F 中每个 x，x 的约化特征多项式是 \(\mathrm{X}^{2}-\mathrm{T}_{\mathrm{F}}(x)\mathrm{X}+\mathrm{N}_{\mathrm{F}}(x)\)。特别地，\(\mathrm{T}_{\mathrm{F}}(x)\) 是 x 的约化迹，\(\mathrm{N}_{\mathrm{F}}(x)\) 是它的约化范数。

(i)⇒(ii)：若代数 F 是中心单的，则由 VIII, 第 361 页，命题 1 与约化迹的定义（VIII, 第 340 页，定义 2）推出，\(T_{F}\) 是它的约化迹；断言 (ii) 随之得出（VIII, 第 343 页，命题 5）。

\begin{enumerate}
\def\labelenumi{(\roman{enumi})}
\setcounter{enumi}{1}
\tightlist
\item
  \(\Leftrightarrow\) (iii)：设 \((e_i)_{1\leqslant i\leqslant 4}\) 是 F 的一个型为 \((\alpha ,\beta ,\gamma)\) 的基（III, §2, No.~5, 第 445 页）。矩阵 \((\mathrm{T_F}(e_i e_j))\) 等于
\end{enumerate}

\[
\left( \begin{array}{c c c c} 2 & \beta & 0 & 0 \\ \beta & 2 \alpha + \beta^ {2} & 0 & 0 \\ 0 & 0 & 2 \gamma & \beta \gamma \\ 0 & 0 & \beta \gamma & - 2 \alpha \gamma \end{array} \right).
\]

它的行列式是 \(-\gamma^{2}(4\alpha + \beta^{2})^{2}\)。性质 (ii) 与 (iii) 的等价性由 V, §8, No.~2, 第 49 页，引理 1 得出。

(iii)⇒(i)：设 \((4\alpha + \beta^{2})\gamma \neq 0\)。则我们有 \(\gamma \neq 0\)，且若 K 的特征为 2，则我们有 \(\beta \neq 0\)。由命题 2，代数 F 是中心的。设 x 是 F 的雅各布森根的一个元素。对每个 \(y \in F\)，xy 是幂零的，故 \(T_{F}(xy) = 0\)（VIII, 第 362 页，注 3）。由于 (ii) 等价于 (iii)，我们有 x = 0。这证明 F 是一个半单 K-代数。由于它的中心是 K，它是单的。

最后的断言由 VIII, 第 361 页，命题 1 与约化特征多项式的定义（VIII, 第 340 页，定义 1）得出。

把 K 的特征记作 p。~由命题 3，若 \(p \neq 2\)，则 K 上每个型为 \((\alpha, 0, \gamma)\)（其中 \(\alpha\) 与 \(\gamma\) 属于 \(K^{*}\)）的四元数代数都是中心单的。若 p = 2，则每个型为 \((\alpha, 1, \gamma)\)（其中 \(\alpha \in K\)，\(\gamma \in K^{*}\)）的四元数代数都是中心单的。反之，我们有下列结果。

命题 4. —— 设 A 是 K 上一个次数为 4 的中心单代数。把 K 的特征记作 p。

\begin{enumerate}
\def\labelenumi{\alph{enumi})}
\item
  若 \(p \neq 2\)，则存在非零元素 \(\alpha\) 与 \(\gamma\)（属于 \(K\)），使得代数 \(A\) 同构于型为 \((\alpha, 0, \gamma)\) 的四元数代数。
\item
  若 \(p = 2\)，则存在一个元素 \(\alpha\)（属于 \(K\)）与一个元素 \(\gamma\)（属于 \(K^*\)），使得代数 \(A\) 同构于型为 \((\alpha, 1, \gamma)\) 的四元数代数。
\end{enumerate}

由韦德伯恩定理（VIII, 第 120 页，定理 1），存在一个整数 \(r \geqslant 1\) 与一个以 K 为中心的体 D，使得 A 同构于 \(\mathbf{M}_r(\mathrm{D})\)。于是我们有 \(r^2[\mathrm{D}:\mathrm{K}] = [\mathrm{A}:\mathrm{K}] = 4\)。若 \(r = 2\)，则 A 同构于 \(\mathbf{M}_2(\mathrm{K})\)，命题 4 由 VIII, 第 362 页的例得出。否则，我们有 \(r = 1\)，且 A 是一个以 K 为中心的体。于是它有一个是 K 的可分扩张的极大交换子域 E；由于

A 在 K 上的次数为 4，故扩张 E 在 K 上的次数为 2（VIII, 第 265 页，推论 2）。因此它是二次的（III, §2, No.~3, 第 439 页）。设 s 是 E 的共轭（III, §2, No.~3, 第 440 页）。由斯科伦--诺特定理（VIII, 第 256 页，推论 1），存在 A 的一个可逆元素 j，使得我们有 \(s(x) = jxj^{-1}\)（对 E 中每个 x）。体 E 在 K 上是可分的，故我们有 \(s \neq Id_{E}\)，从而 \(j \notin E\)。由于 A 是 K 上一个维数为 4 的向量空间，它是 E 上一个维数为 2 的左向量空间，于是我们有 \(A = E \oplus Ej\)。我们有 \(s^{2} = Id_{E}\)，故 A 的元素 \(j^{2}\) 属于 A 的中心；因此存在一个元素 \(\gamma\)（属于 \(K^{*}\)）使得 \(j^{2} = \gamma\)。

当 \(p \neq 2\) 时，存在 E 的一个元素 i 与一个元素 \(\alpha \in K^{*}\)，使得 \(\mathrm{E} = \mathrm{K}(i)\) 且 \(i^{2} = \alpha\)（V, §11, No.~9, 第 93 页，例 3）；在这种情形，A 同构于型为 \((\alpha, 0, \gamma)\) 的四元数代数。当 p = 2 时，存在 E 的一个元素 i 与 K 的一个元素 \(\alpha\)，使得 \(\mathrm{E} = \mathrm{K}(i)\) 且 \(i^{2} = i + \alpha\)（V, §11, No.~9, 第 93 页，例 2），因而 A 同构于型为 \((\alpha, 1, \gamma)\) 的四元数代数。

推论 1. —— 设 A 是一个有限次数 \textgreater{} 1 的中心单 K-代数，其元素在 K 上都是次数 \(\leqslant 2\) 的代数元。则 A 同构于 K 上的一个四元数代数。

若 K 有限，则代数 A 同构于一个矩阵代数 \(\mathbf{M}_{n}(\mathrm{K})\)（VIII, 第 357 页，推论 2），因而包含 K 上次数为 n 的元素；假设蕴含 n = 2，从而得到该情形下的结果（VIII, 第 362 页，例）。设域 K 是无限的。设 L 是 A 的一个极大平展子代数。由 V, §7, No.~4, 第 41 页，命题 7，存在 A 的一个元素 x，使得 K-代数 L 等于 \(K[x]\)，故由假设其次数 \(\leqslant 2\)。由于我们有 \([A : K] = [L : K]^{2}\)（VIII, 第 264 页，命题 4 与第 262 页，命题 3），我们断定 \([A : K] = 4\)。于是推论 1 由命题 4 得出。

推论 2. —— 设 \((\mathrm{E}, s)\) 是 K 上的一个凯莱代数，使得 K-代数 E 是有限次数 \textgreater{} 1 的中心单代数。则 E 同构于 K 上的一个四元数代数。

每个元素 \(u\)（属于 \(\mathrm{E}\)）满足 \(u^2 - \mathrm{T}_{\mathrm{E}}(u)u + \mathrm{N}_{\mathrm{E}}(u) = 0\)，故 K-代数 \(\mathrm{E}\) 同构于一个四元数代数（推论 1）。

